\documentclass[10pt,a4paper]{amsart}
\usepackage[margin=19mm]{geometry}
\usepackage{amsmath,amssymb,mathtools}
\usepackage[scr=rsfs,cal=euler]{mathalfa}
\usepackage{xfrac}
\usepackage[unicode,hidelinks,bookmarks=false]{hyperref}
\newcommand{\ie}{i.e.\ }
\newcommand{\cf}{cf.\ }
\newcommand{\resp}{respectively\ }
\newcommand{\Subsection}{\subsection}
\providecommand{\nobreakdash}{\nobreak\hbox{-}}
\setlength{\emergencystretch}{2em}
\multlinegap0pt
\usepackage{xcolor}
\newtheorem{prop}{Proposition}
\newtheorem{thm}{Theorem}
\usepackage{cleveref}
\crefname{prop}{Proposition}{Propositions}
\crefname{thm}{Theorem}{Theorems}
\newcommand{\N}{\mathbb N}
\newcommand{\mK}{\mathcal{K}}
\newcommand{\obp}{\otimes^{\gamma}}
\newcommand{\ot}{\otimes}
\title{On strong Arens irregularity of projective tensor product of Hilbert-Schmidt space}
\author{Ved Prakash Gupta, Lav Kumar Singh}
\date{Source: arXiv:2108.13946v2 (CC BY 4.0)}
\begin{document}
\maketitle

\noindent\emph{Source and licence.} On strong Arens irregularity of projective tensor product of Hilbert-Schmidt space. Version arXiv:2108.13946v2 (CC BY 4.0), distributed by arXiv under the Creative Commons Attribution 4.0 International licence, http://creativecommons.org/licenses/by/4.0/, as recorded in the arXiv licence metadata for this exact version and shown on the abstract page https://arxiv.org/abs/2108.13946v2. Full licence notice: \texttt{licenses/arxiv-2108.13946-CC-BY-4.0.txt}.\par\smallskip

\noindent\textit{Editorial scope.} Counted unit: the article's thm ann-box (source0.tex lines 662--666). The article's complete original proof of it is delivered with it (source0.tex lines 667--733, one \texttt{proof} environment of 67 lines). No mathematics is written, repaired or re-derived: the statement and the proof are the article's own lines, in the article's own notation, and no retained line is substituted or re-flowed.  Retained uncounted as the source-local context the proof needs: lines 250--258; lines 493--498.  Nothing was omitted from any retained span, so the package carries no omission record.  No retained line contains a citation, so the wrapper carries no bibliography and no bibliography entry.  No retained line contains a figure environment, an image-inclusion command or drawing code, so no image is delivered and the package declares no asset dependency.  Editorial formatting only, in addition: an AMS \texttt{amsart} preamble that restates the article's own declarations and macros used by the retained lines, namely the environments \texttt{prop}, \texttt{thm} on separate counters, together with the standard packages those lines need.  The retained environments print in this document as Proposition 1, Theorem 1 in order of appearance, whereas the article prints the same items with the numbering of its own full document.  The complete native source of the article as distributed in the arXiv e-print for this exact version is bundled unchanged as \texttt{sources/115526/source0.tex}.
\begin{prop}\label{Th:2.1}
	Let $\mathcal H_1,\mathcal H_2$ and $\mathcal H_3$ be Hilbert
        spaces. If $R\in B(\mathcal H_2)$, $S\in S_p(\mathcal
        H_1,\mathcal H_2)$ and $T\in B( \mathcal H_1)$,
        then
        \[
        \|RST\|_p\leq \|R\| \|S\|_p \|T\|.
        \]
\end{prop}

\begin{equation} \label{eqn:3.1}
(\phi_{\mathcal U}\diamond
  \psi)(f)=\psi\left(\overset{w^*}{\lim_{\mathcal U}} f_{ e_{nn}\otimes
    e_{nn}}\right) \mathrm{ and }\ (\phi_{\mathcal U}\Box \psi)(f) =
  \lim_{\mathcal U}\psi(f_{ e_{nn} \otimes e_{nn}})
\end{equation}

\begin{thm}\label{ann-box}
	For any non-principal ultrafilter $\mathcal U$ on $\mathbb
        N$, $$\mathcal{A}^r_{(A^{**},\Box)}\left(\{\phi_{\mathcal
          U}\}\right)=A^{**}.$$
	\end{thm}

\begin{proof}
Let $\psi\in (\mK \obp \mK)^{**}$ and $f\in (\mathcal K\obp\mathcal
K)^{*}$. Then, from \Cref{eqn:3.1}, we know that $(\phi_{\mathcal U}
\Box \psi)(f) = \lim_{\mathcal{ U}} \psi(f_{e_{nn} \ot e_{nn}})$.
Since $\mathcal{U}$ contains the co-finite ultrafilter, it can be
easily seen that $\lim_{\mathcal{ U}} \psi(f_{e_{nn} \ot e_{nn}})$ equals
$\lim_{n\to \infty} \psi(f_{e_{nn} \ot e_{nn}})$ if the latter exists. So, in
order to show that $ (\phi_{\mathcal U}\Box \psi)(f)=0 $, it suffices
to show that $\lim_{n\to \infty}\psi(f_{ e_{nn} \otimes e_{nn}}) =0$.

To prove this, let us assume on the contrary that the sequence
$\{\psi(f_{e_{nn} \otimes e_{nn}})\}$ does not converge to $0$ as $n$
tends to infinity. Then, there exists an $\epsilon>0$ for which we can choose an
increasing sequence $\{n_t\}_{t=1}^\infty$ of natural numbers such
that $|\psi(f_{ e_{n_tn_t} \otimes e_{n_tn_t}})|>\epsilon$ for all $t
\geq 1$.  Let us define $g_{t}\in A^*$ for each natural number $t$ as
$g_{t}=c_{t}\cdot f_{ e_{n_tn_t}\otimes e_{n_tn_t}},$ where
$c_t=\frac{\overline{\psi(f_{ e_{n_tn_t} \otimes
      e_{n_tn_t}})}}{\left\vert\psi(f_{e_{n_tn_t} \otimes
    e_{n_tn_t}})\right\vert}$. Now, for each $N \in \N$, let
$h_N:=\sum_{t=1}^Ng_{t}$. Then,
\[
	|h_N(u)| = \left\vert f\left (\sum_{t=1}^N c_t(
        e_{n_tn_t}\otimes e_{n_tn_t})\cdot u\right)\right\vert \leq
        ||f||\,
        \left\vert\left\vert\sum_{t=1}^N(c_tu_{n_t})\right\vert\right\vert_{\gamma}
\]
for all $u \in \mK\obp \mK$ and $N \in \N$, where $u_{n_t}=(e_{n_tn_t}\otimes e_{n_tn_t})\cdot u$.

For each $N \in \N$, consider the linear operator $Q_N:\mathcal K\to
\mathcal K$ satisfying $Q_N(e_{n_tj})=c_t e_{n_tj}$ for $t=1,2,..,N$,
$j\in \mathbb N$ and $Q_N(e_{ij})=0$ otherwise. Clearly,
$||Q_N||=1$. Then, one can easily verify, through the actions on the
orthonormal basis $\{e_{ij}\}$, that
\[
\theta\left(\sum_{t=1}^Nc_tu_{n_t}\right)=Q_N^*\theta(u)Q_N
\]
for all $u \in \mK \obp \mK$ and $N \in \N$. Thus, from \Cref{Th:2.1},
we obtain
\[
\left\|\theta\left(\sum_{t=1}^Nc_tu_{n_t}\right)\right\|_1=\left\|\sum_{t=1}^Nc_tu_{n_t}\right\|_\gamma\leq
||u||_\gamma
\]
for all $u \in \mK \obp \mK$ and $N \in \N$. In particular,
\[
  |h_N(u)|\leq |f||||u||_\gamma
  \]
  for all $u \in \mK\obp \mK$ and $N \in \N$.  Thus, $||h_N||\leq
  ||f||$ for each natural number $N$.

  On the other hand,
  \[
  \psi(h_N)=\sum_{t=1}^N\psi(g_t)=\sum_{t=1}^N|\psi(f_{
    e_{n_tn_t}\otimes e_{n_tn_t}})| \text{ for all } N \in \N.
  \]
  Thus, $\{\psi(h_N): N \in \N\}$ is an unbounded sequence. But this
  is absurd because
  \[
  |\psi(h_N)|\leq \|\psi \|\|h_N\| \leq \|\psi\| \|f\|
  \]
  for all $N \in \N$. Hence, our assumption was wrong and we must have $\lim_{n\to
    \infty}\psi(f_{ e_{nn}\otimes e_{nn}})=0$.

  Since $\psi$ and $f$ were arbitrary, it follows that
  $\mathcal{A}^r_{(A^{**},\Box)}\left(\{\phi_{\mathcal
    U}\}\right)=A^{**}.$
	\end{proof}

\end{document}
